\section{El campo de supresiones y sus derivadas}
Sea $E$ el espacio de triples $(h,t_R,t_C)$ de masas iguales, con norma suma de TV.
Para $U$ RC compatible se pone $T_U=\sum_{e\in U}(1,\delta_{r_e},\delta_{c_e})$.
La coordenada analítica es $T$, de escala de conteo; $\theta=T/j$.
Denotemos por $\mu_0,\nu_0$ las medidas de masa $1/j$ por índice. Balanceamos $z=jD$
en las medidas $(\mu_0-\theta_R)/(1-\theta_0)$ y $(\nu_0-\theta_C)/(1-\theta_0)$.
Así $z'=z\exp(s_r+t_c)$ tiene ambos márgenes uno. El gauge es la media original
cero de $t$. Basta imponer las filas y las columnas proyectadas: el residuo constante
en las columnas desaparece al integrar ambos sistemas con las nuevas medidas.

En el centro el complemento de Schur en columnas centradas es $I-D^tD$.
La brecha acota el inverso en norma dos normalizada. Las ecuaciones
por fila y columna convierten la cota dos en cota en supremo, pues el núcleo $jD$ es acotado.
En una bola fija las segundas derivadas del sistema son uniformes; su variación al cambiar
las medidas es $O(\|\theta\|_{\rm TV})$ en la norma de supremo. Por la iteración con
inversa inicial congelada existe una rama holomorfa acotada con radio de datos fijo.
Se conservan las ecuaciones auxiliares de los índices de masa cero. En los puntos reales
el paso a la representación con raíces de las medidas cambia el operador en
$O(\sqrt{h/j})$ en Hilbert--Schmidt, y la rama conserva márgenes y brecha (U02, 27--38).

En la bola compleja el operador $O$ de U02 es el núcleo cruzado $z'-1$.
En el centro real, en $|yw|r'<1$ con margen, el resolvente en supremo se acota
escribiéndolo como $I-ywO+(yw)^2 O(I+ywO)^{-1}O$. Los dos núcleos exteriores
convierten entre normas dos y supremo. Una perturbación pequeña en TV y en el núcleo
preserva la inversa por Neumann; aplicarla a las columnas acota el propagador puntual.
Los sumandos de Fredholm se acotan por $C^l l^{l/2}/l!$ por Hadamard.
La serie y la variación por telescopaje son uniformes. En el centro la separación
del determinante de cero es uniforme por la brecha y la suma de cuadrados singulares.
Se obtiene separación y rama acotada en una bola de datos reducida. Esto justifica
las cotas complejas usadas en la parte anterior con medidas deformadas. La inversión de
coincidencias de U03 se prolonga con una integral por bloque, según su definición.

\paragraph{Normalización exacta.}
Pongamos $A_0(\theta)=-\int s\,d(\mu_0-\theta_R)-\int t\,d(\nu_0-\theta_C)$.
En el punto físico, en los supervivientes,
$D_{-U}=(m/j)\operatorname{diag}(e^{-s})D'\operatorname{diag}(e^{-t})$.
Por ello, para $D_e>0$ en $U$ y $p_U=\mathbb P_{\rm torres}(U\subset\sigma)$,
\begin{align*}
 \log(p_U/\textstyle\prod_{e\in U}D_e)
 &=\log(m!/j!)+h\log j+j A_0(\theta)\\
 &\quad+\mathcal L_{Q_*}(1/m;z')-\mathcal L_{Q_*}(1/j;z)+O(e^{-cj}),
\end{align*}
donde $\mathcal L_{Q_*}=\log H_{Q_*}$ usa las correspondientes medidas balanceadas.
El signo y el factor $j^h$ se obtienen multiplicando una vez por cada margen superviviente.
$A_0(0)=0$ y la acción es invariante al gauge. Conservamos el mismo entero $Q_*$.
Sea $b_e=D_e\exp(A'_0(0)v_e)$, $v_e=(1,\delta_{r_e},\delta_{c_e})$.
Llamamos $L_s(T)$ al miembro derecho sin el error y restando $A'_0(0)T$.
En la bola $\|T\|\le a_0 j$ la expresión factorial se prolonga por $\log\Gamma$;
integrar $\psi(j+1-h)-\log j=O(1)$ allí da una cota $Cj$.
Elegimos $\varepsilon$ de modo que $|1/(j-h)|<c_0/(2Q_*)$ en la bola.
La parte factorial tiene derivada en cero $\log j-\psi(j+1)=O(1/j)$.
La acción centrada tiene derivada cero. En $\mathcal L_{Q_*}$ la derivada en tamaño
cuesta $CQ_*/j^2$ y en datos $C/j$, por los dos radios distintos. Resulta
\[
 L_s(0)=0,\quad\sup|L_s|\le Cj,\quad
 \|DL_s(0)\|\le C/j,\quad \sup_{\|T\|\le a_0 j/2}\|D^2 L_s(T)\|\le C/j.
\]
La última desigualdad es Cauchy y polarización. Las derivadas pertenecen al campo
holomorfo. No se deriva el error entre dos valores físicos discretos.
Para $h\le a_* j$ con $3a_*<a_0/2$ resulta
$p_U=(\prod b_e)e^{L_s(T_U)}(1+O(e^{-cj}))$. Para todo $h$ se tiene
$p_U\le C_0^h\prod b_e$: hasta $a_* j$ se usa el campo, y después
$\operatorname{per}(D_{-U})\le1$, $\operatorname{per}D\ge j!/j^j\ge e^{-j}$
y $j\le h/a_*$. Ésta es la reserva independiente del corte de Taylor posterior.

El orden común se fija con el padre: $Q_*=\lfloor\varepsilon j\rfloor$, incluso en el menor.
